\documentclass[11pt]{amsart}

\usepackage[margin=1in]{geometry}
\usepackage{amsmath,amssymb,amsthm,mathtools}
\usepackage{hyperref}

\newtheorem{theorem}{Theorem}[section]
\newtheorem{proposition}[theorem]{Proposition}
\newtheorem{lemma}[theorem]{Lemma}
\newtheorem{corollary}[theorem]{Corollary}

\theoremstyle{remark}
\newtheorem{remark}[theorem]{Remark}

\newcommand{\C}{\mathbb{C}}
\newcommand{\N}{\mathbb{N}}
\newcommand{\R}{\mathbb{R}}
\newcommand{\Z}{\mathbb{Z}}
\newcommand{\abs}[1]{\lvert #1 \rvert}
\newcommand{\norm}[1]{\lVert #1 \rVert}
\newcommand{\ip}[2]{\langle #1,#2\rangle}
\newcommand{\dd}{\,\mathrm{d}}
\newcommand{\ee}{\mathrm{e}}
\newcommand{\ii}{\mathrm{i}}
\newcommand{\rhoFn}{\rho}
\DeclareMathOperator{\dist}{dist}

\title{Local Stable Phase Retrieval in the First Hermite Model Space}

\begin{document}

\begin{abstract}
In the analytic Fock space, local stable phase retrieval at the constant
function is proved by combining an abstract orthogonal reduction with circle
estimates for positive-frequency trigonometric polynomials and a radial
localisation argument. This note records the corresponding argument in the
first true polyanalytic model space associated with the first Hermite window.
The new input is an annulus localisation estimate for the basis functions of
that model space. Relative to the earlier proof sketch, we make the dependence
on the block width explicit in the annulus-circle estimate, define the leakage
parameter quantitatively, repair the signed squaring step in the assembly
argument, and state the orthogonal reduction only in the phase-normalised form
that is supplied by the verified analytic proof. The final theorem is therefore
formulated entirely in the model space $H_1$, where it yields local
phase-optimised stability near the distinguished vector $\Phi_0=\overline z$.
\end{abstract}

\maketitle

\section{Introduction}

Let
\[
L^2_\gamma(\C)
:=
L^2\!\Bigl(\C,\frac{\ee^{-\abs{z}^2}}{\pi}\,\dd m(z)\Bigr),
\]
with inner product
\[
\ip{F}{G}
:=
\frac1\pi \int_\C F(z)\overline{G(z)}\,\ee^{-\abs{z}^2}\,\dd m(z),
\qquad
\norm{F}^2 := \ip{F}{F}.
\]

For the first Hermite window, the relevant model space is the first true
polyanalytic level. We work with the functions
\[
\Phi_0(z) := \overline z,
\qquad
\Phi_n(z) := \frac{z^{n-1}}{\sqrt{n!}}\bigl(\abs{z}^2 - n\bigr)
\quad (n \ge 1),
\]
and define
\[
H_1 := \overline{\operatorname{span}}\{\Phi_n : n \ge 0\} \subset L^2_\gamma(\C).
\]
Thus $H_1$ is a Hilbert space by construction, and below we will see that
$(\Phi_n)_{n\ge0}$ is an orthonormal basis of $H_1$.

Our goal is to prove the Hermite analogue of the analytic orthogonal coercivity
estimate:
\[
\norm{G}
\lesssim
\norm{\abs{\Phi_0 + G} - \abs{\Phi_0}}
\qquad
\bigl(G \in H_1 \cap \Phi_0^\perp\bigr).
\]
Once this is done, the phase-normalised orthogonal reduction from the verified
analytic proof, together with an elementary phase-alignment lemma, yields local
phase-optimised stability near $\Phi_0$ inside $H_1$.

The overall architecture is the same as in the analytic Fock-space proof:
\begin{enumerate}
\item identify an orthonormal basis adapted to polar coordinates;
\item reduce the modulus defect on each circle to the nonlinear quantity
\[
\rhoFn(w) := \abs{1+w} - 1
\]
applied to a positive-frequency trigonometric polynomial;
\item localise square frequency blocks to spatial annuli;
\item combine a local-circle estimate, a high-frequency estimate, and a
leakage argument.
\end{enumerate}
The only genuinely new estimate is the annulus localisation for the basis
functions $\Phi_n$.

The theorem proved here is a theorem in the model space $H_1$. We do not
reproduce in this note the identification between $H_1$ and the image of the
short-time Fourier transform with the first Hermite window, so we do not state
here a separate signal-space corollary.

Throughout, we use
\[
\N = \{0,1,2,\dots\}.
\]

\section{Main theorem}

We use the following ingredients from the verified analytic Fock-space proof as
black boxes. Here we keep the same signed defect $\rhoFn(w)=\abs{1+w}-1$ as in
that argument, so the quoted circle estimates apply without any change of
notation.

\begin{proposition}[Local circle estimate]\label{prop:local-circle}
Let $E \subset \Z_{>0}$ be finite with $\#E \le L$, and let
\[
P(\zeta) = \sum_{n\in E} c_n \zeta^n
\]
be a trigonometric polynomial on $S^1$ with only strictly positive
frequencies. Then
\[
\int_{S^1} \abs{P(\zeta)}^2\,\dd \sigma(\zeta)
\le
144\,L \int_{S^1} \rhoFn\bigl(P(\zeta)\bigr)^2\,\dd \sigma(\zeta).
\]
\end{proposition}

\begin{proposition}[High-frequency circle estimate]\label{prop:high-circle}
Let $N \ge 1$, and let
\[
P(\zeta) = \sum_{n=N}^{N+L-1} c_n \zeta^n
\]
have only strictly positive frequencies. If
\[
1343\,L^2 \le N^2,
\]
then
\[
\int_{S^1} \abs{P(\zeta)}^2\,\dd \sigma(\zeta)
\le
32 \int_{S^1} \rhoFn\bigl(P(\zeta)\bigr)^2\,\dd \sigma(\zeta).
\]
\end{proposition}

\begin{proposition}[Phase-normalised orthogonal reduction]\label{prop:orth-red}
Let $H$ be a complex Hilbert space continuously embedded in an $L^2$-space,
and let $f_0 \in H$ satisfy $\norm{f_0}=1$. Assume there exists $C>0$ such
that
\[
\norm{g}
\le
C\,\norm{\abs{f_0 + g} - \abs{f_0}}
\qquad
\text{for all } g \in H \cap f_0^\perp.
\]
Then there exist $\delta,\widetilde M>0$ such that for every $h \in H$ with
$\norm{h}\le\delta$ and $\ip{h}{f_0}\in\R$, one has
\[
\norm{h}
\le
\widetilde M\,\norm{\abs{f_0 + h} - \abs{f_0}}.
\]
\end{proposition}

\begin{lemma}[Phase alignment]\label{lem:phase-alignment}
Let $H$ be a complex Hilbert space, let $f_0 \in H$ with $\norm{f_0}=1$, and
let $u \in H$. Then there exists $w \in \C$, $\abs{w}=1$, such that
\[
h := w(f_0 + u) - f_0
\]
satisfies $\ip{h}{f_0} \in \R$ and $\norm{h} \le \norm{u}$.
\end{lemma}

\begin{proof}
Set
\[
\beta := 1 + \ip{u}{f_0}.
\]
If $\beta = 0$, take $w=1$. Then
\[
\ip{h}{f_0} = \ip{u}{f_0} = -1 \in \R
\]
and $\norm{h}=\norm{u}$.

Assume now $\beta \ne 0$ and choose
\[
w := \frac{\overline{\beta}}{\abs{\beta}}.
\]
Then
\[
\ip{h}{f_0}
=
w\,\ip{f_0 + u}{f_0} - 1
=
w\beta - 1
=
\abs{\beta} - 1
\in \R.
\]
Moreover,
\begin{align*}
\norm{h}^2
&=
\norm{w(f_0 + u) - f_0}^2 \\
&=
\norm{f_0 + u}^2 + \norm{f_0}^2 - 2\,\Re\!\bigl(w\,\ip{f_0 + u}{f_0}\bigr) \\
&=
\bigl(1 + 2\Re\ip{u}{f_0} + \norm{u}^2\bigr) + 1 - 2\abs{\beta} \\
&=
\norm{u}^2 + 2\bigl(1 + \Re\ip{u}{f_0} - \abs{\beta}\bigr)
\le
\norm{u}^2,
\end{align*}
because $\abs{\beta}\ge\Re(\beta)=1+\Re\ip{u}{f_0}$.
\end{proof}

\begin{theorem}\label{thm:main-orthog}
There exists a constant $C>0$ such that every
$G \in H_1 \cap \Phi_0^\perp$ satisfies
\[
\norm{G}
\le
C\,\norm{\abs{\Phi_0 + G} - \abs{\Phi_0}}.
\]
\end{theorem}

\begin{corollary}\label{cor:main-local}
There exist $\delta,\widetilde M>0$ such that every $H \in H_1$ with
$\norm{H}\le\delta$ admits a phase $w \in \C$, $\abs{w}=1$, with
\[
\norm{w(\Phi_0 + H) - \Phi_0}
\le
\widetilde M\,\norm{\abs{\Phi_0 + H} - \abs{\Phi_0}}.
\]
\end{corollary}

\begin{proof}
Apply Lemma~\ref{lem:phase-alignment} with $f_0=\Phi_0$ and $u=H$ to obtain
$w$ such that
\[
h := w(\Phi_0 + H) - \Phi_0
\]
satisfies $\ip{h}{\Phi_0}\in\R$ and $\norm{h}\le\norm{H}$. By
Lemma~\ref{lem:orthonormal}, $\norm{\Phi_0}=1$. If
$\delta>0$ is chosen as in Proposition~\ref{prop:orth-red} for the constant
from Theorem~\ref{thm:main-orthog}, then $\norm{h}\le\delta$ whenever
$\norm{H}\le\delta$. Moreover,
\[
\abs{\Phi_0 + h} - \abs{\Phi_0}
=
\abs{w(\Phi_0 + H)} - \abs{\Phi_0}
=
\abs{\Phi_0 + H} - \abs{\Phi_0}.
\]
Proposition~\ref{prop:orth-red} therefore yields the claim.
\end{proof}

\section{The first true-level basis}

For $n \ge 0$, let
\[
e_n(z) := \frac{z^n}{\sqrt{n!}}.
\]

Here and below, for $z=x+iy$ we write
\[
\partial_z := \frac12\bigl(\partial_x - i\partial_y\bigr)
\]
for the standard Wirtinger derivative.

\begin{lemma}\label{lem:basis-operator}
For every $n \ge 0$,
\[
\Phi_n = (\overline z - \partial_z)e_n.
\]
\end{lemma}

\begin{proof}
For $n=0$, this is $\Phi_0 = \overline z$. For $n \ge 1$,
\[
(\overline z - \partial_z)e_n(z)
=
\frac{\overline z z^n}{\sqrt{n!}} - \frac{n z^{n-1}}{\sqrt{n!}}
=
\frac{z^{n-1}}{\sqrt{n!}}\bigl(\abs{z}^2 - n\bigr)
=
\Phi_n(z).
\qedhere
\]
\end{proof}

\begin{lemma}\label{lem:polar-basis}
If $z=re^{it}$, then
\[
\Phi_0(re^{it}) = re^{-it},
\]
and for every $n \ge 1$,
\[
\Phi_n(re^{it})
=
\frac{r^{n-1}(r^2 - n)}{\sqrt{n!}}\,e^{i(n-1)t}.
\]
\end{lemma}

\begin{proof}
The first identity is immediate from $\Phi_0(z)=\overline z$. For $n \ge 1$,
\[
\Phi_n(re^{it})
=
\frac{(re^{it})^{n-1}}{\sqrt{n!}}(r^2-n)
=
\frac{r^{n-1}(r^2 - n)}{\sqrt{n!}}\,e^{i(n-1)t}.
\qedhere
\]
\end{proof}

\begin{lemma}\label{lem:orthonormal}
The family $(\Phi_n)_{n\ge0}$ is orthonormal in $L^2_\gamma(\C)$.
\end{lemma}

\begin{proof}
By Lemma~\ref{lem:polar-basis}, $\Phi_0(re^{it})$ has angular frequency $-1$,
while $\Phi_n(re^{it})$ has angular frequency $n-1$ for $n \ge 1$. Hence
distinct basis vectors are orthogonal on each circle and therefore in
$L^2_\gamma(\C)$.

For $n=0$,
\[
\norm{\Phi_0}^2
=
\frac1\pi \int_\C \abs{z}^2 \ee^{-\abs{z}^2}\,\dd m(z)
=
2\int_0^\infty r^3 \ee^{-r^2}\,\dd r
=
\int_0^\infty u\,\ee^{-u}\,\dd u
=
1.
\]

For $n \ge 1$, Lemma~\ref{lem:polar-basis} gives
\[
\norm{\Phi_n}^2
=
\frac{2}{n!}\int_0^\infty r^{2n-1}(r^2-n)^2 \ee^{-r^2}\,\dd r
=
\frac{1}{n!}\int_0^\infty u^{n-1}(u-n)^2 \ee^{-u}\,\dd u.
\]
Expanding $(u-n)^2$ and using the Gamma integrals,
\[
\norm{\Phi_n}^2
=
\frac{(n+1)! - 2n\,n! + n^2 (n-1)!}{n!}
=
1.
\qedhere
\]
\end{proof}

\begin{corollary}\label{cor:expansion}
Every $G \in H_1$ has a unique expansion
\[
G = \sum_{n=0}^\infty a_n \Phi_n
\qquad \text{in } H_1,
\]
with
\[
\norm{G}^2 = \sum_{n=0}^\infty \abs{a_n}^2.
\]
Moreover, $G \in \Phi_0^\perp$ if and only if $a_0=0$.
\end{corollary}

\begin{proof}
This is immediate from the definition of $H_1$ as the closed span of the
orthonormal family $(\Phi_n)_{n\ge0}$.
\end{proof}

\section{Reduction to the circle}

The crucial observation is that after dividing by the distinguished vector
$\Phi_0$, the perturbation still has only strictly positive angular
frequencies.

\begin{lemma}\label{lem:circle-reduction}
Let
\[
G = \sum_{n=1}^D a_n \Phi_n.
\]
For $r>0$, define
\[
P_r(e^{it})
:=
\sum_{n=1}^D a_n c_n(r)e^{int},
\qquad
c_n(r) := \frac{r^{n-2}(r^2-n)}{\sqrt{n!}}.
\]
Then
\[
G(re^{it}) = r e^{-it} P_r(e^{it}),
\qquad
\Phi_0(re^{it}) + G(re^{it}) = r e^{-it}\bigl(1 + P_r(e^{it})\bigr).
\]
Consequently,
\[
\bigl(\abs{\Phi_0 + G} - \abs{\Phi_0}\bigr)(re^{it})
=
r\,\rhoFn\bigl(P_r(e^{it})\bigr).
\]
\end{lemma}

\begin{proof}
By Lemma~\ref{lem:polar-basis},
\[
\Phi_n(re^{it})
=
r e^{-it}\frac{r^{n-2}(r^2-n)}{\sqrt{n!}}\,e^{int}
\qquad (n \ge 1).
\]
Summing in $n$ gives the identity for $G$. The formula for $\Phi_0+G$ follows
from $\Phi_0(re^{it})=re^{-it}$. Since
\[
\abs{\Phi_0(re^{it}) + G(re^{it})}
=
r\abs{1 + P_r(e^{it})}
\quad \text{and} \quad
\abs{\Phi_0(re^{it})}=r,
\]
subtracting gives the last claim.
\end{proof}

\section{Annulus localisation for the basis functions}

We first isolate the elementary ``polynomial absorbs Gaussian'' estimate used
twice in the proof of annulus localisation.

\begin{lemma}\label{lem:poly-gauss}
Let $a>0$ and let $k \in \N$. There exists $C_{a,k}>0$ such that
\[
\bigl(1+x^k\bigr)\ee^{-a x^2}
\le
C_{a,k}\,\ee^{-a x^2/2}
\qquad
(x \ge 0).
\]
\end{lemma}

\begin{proof}
The function $x \mapsto (1+x^k)\ee^{-a x^2/2}$ is continuous on $[0,\infty)$
and tends to $0$ as $x\to\infty$, so it is bounded.
\end{proof}

For $j \ge 0$, let
\[
A_j := \{z \in \C : j \le \abs{z} < j+1\}.
\]

\begin{lemma}\label{lem:annulus-basis}
There exist absolute constants $C,c>0$ such that for every $n \ge 1$ and
every $j \ge 0$,
\[
\frac1\pi \int_{A_j} \abs{\Phi_n(z)}^2 \ee^{-\abs{z}^2}\,\dd m(z)
\le
C \exp\!\Bigl(-c\bigl(\abs{j-\sqrt n}-3\bigr)_+^2\Bigr).
\]
\end{lemma}

\begin{proof}
Set
\[
I_{n,j}
:=
\frac1\pi \int_{A_j} \abs{\Phi_n(z)}^2 \ee^{-\abs{z}^2}\,\dd m(z).
\]
By Lemma~\ref{lem:polar-basis} and polar coordinates,
\[
I_{n,j}
=
\frac{2}{n!}\int_j^{j+1} r^{2n-1}(r^2-n)^2 \ee^{-r^2}\,\dd r
=
\frac1n \int_{j^2}^{(j+1)^2} (u-n)^2 q_n(u)\,\dd u,
\]
where
\[
q_n(u) := \frac{u^{n-1}\ee^{-u}}{(n-1)!}
\]
is the $\Gamma(n,1)$ density.

We first bound $q_n$. There exists $C_1>0$ such that for all $n \ge 1$ and
$u \ge 0$,
\[
q_n(u)
\le
\frac{C_1}{\sqrt n}
\exp\!\bigl(-(\sqrt u - \sqrt{n-1})^2\bigr).
\]
For $n=1$, this is immediate. For $n \ge 2$, let $a=n-1$ and set
$u=v^2$, $\psi_a(v)=2a\log v - v^2$. Then
\[
\psi_a''(v) = -\frac{2a}{v^2} - 2 \le -2,
\]
so $\psi_a$ is concave and attains its maximum at $v=\sqrt a$. Hence
\[
\psi_a(v) \le \psi_a(\sqrt a) - (v-\sqrt a)^2.
\]
Exponentiating and dividing by $(n-1)! = a!$, Stirling's lower bound gives the
claim.

Now set
\[
y := \abs{\sqrt u - \sqrt n}.
\]
Since
\[
\abs{u-n} = y(\sqrt u + \sqrt n)
\quad \text{and} \quad
\sqrt u \le y + \sqrt n,
\]
we obtain
\[
\frac{(u-n)^2}{n}
=
\frac{y^2(\sqrt u + \sqrt n)^2}{n}
\le
2y^2 \frac{u+n}{n}
\le
8y^2 + 2y^4
\le
10(1+y^4).
\]
Also,
\[
y
\le
\abs{\sqrt u - \sqrt{n-1}} + \abs{\sqrt n - \sqrt{n-1}}
\le
\abs{\sqrt u - \sqrt{n-1}} + 1.
\]
Therefore, after enlarging constants,
\[
\frac{(u-n)^2}{n}q_n(u)
\le
\frac{C_2}{\sqrt n}
\Bigl(1 + \abs{\sqrt u - \sqrt{n-1}}^4\Bigr)
\exp\!\bigl(-(\sqrt u - \sqrt{n-1})^2\bigr).
\]
Applying Lemma~\ref{lem:poly-gauss} with $a=1$ and $k=4$, we obtain
\[
\frac{(u-n)^2}{n}q_n(u)
\le
\frac{C_3}{\sqrt n}
\exp\!\Bigl(-\frac12(\sqrt u - \sqrt{n-1})^2\Bigr).
\]

Now set
\[
d := \dist\bigl([j,j+1],\sqrt{n-1}\bigr).
\]
For every $u \in [j^2,(j+1)^2]$, one has
$\abs{\sqrt u - \sqrt{n-1}} \ge d$, hence
\[
I_{n,j}
\le
\frac{C_3}{\sqrt n}\,\ee^{-d^2/2}\bigl((j+1)^2 - j^2\bigr)
=
\frac{C_3(2j+1)}{\sqrt n}\,\ee^{-d^2/2}.
\]

We now compare $d$ with $\abs{j-\sqrt n}$. Since
$\abs{\sqrt n - \sqrt{n-1}} \le 1$,
\[
d
\ge
\dist\bigl([j,j+1],\sqrt n\bigr) - 1.
\]
Moreover,
\[
\dist\bigl([j,j+1],\sqrt n\bigr)
\ge
\bigl(\abs{j-\sqrt n}-1\bigr)_+,
\]
so
\[
d \ge \bigl(\abs{j-\sqrt n}-2\bigr)_+.
\]

If $\abs{j-\sqrt n}\le 3$, then the desired right-hand side is bounded below
by a positive constant, while $I_{n,j}\le\norm{\Phi_n}^2=1$ by
Lemma~\ref{lem:orthonormal}, so the claim is automatic after enlarging $C$.

Assume now $\abs{j-\sqrt n}>3$. Then
\[
\frac{2j+1}{\sqrt n}
\le
2 + \frac{2\abs{j-\sqrt n}+1}{\sqrt n}
\le
3 + 2\abs{j-\sqrt n},
\]
so
\[
I_{n,j}
\le
C_4\bigl(1+\abs{j-\sqrt n}\bigr)
\exp\!\Bigl(-\frac12\bigl(\abs{j-\sqrt n}-2\bigr)^2\Bigr).
\]
Applying Lemma~\ref{lem:poly-gauss} once more with $a=\frac12$ and $k=1$,
and then weakening constants slightly, gives
\[
I_{n,j}
\le
C \exp\!\Bigl(-c\bigl(\abs{j-\sqrt n}-3\bigr)_+^2\Bigr).
\qedhere
\]
\end{proof}

\section{Block decomposition}

As in the analytic proof, we divide the frequency axis into square blocks. For
$\ell \ge 0$, define
\[
I_\ell := \{n \in \N : \ell^2 \le n < (\ell+1)^2\}.
\]
Given a finite sum
\[
G = \sum_{n=1}^D a_n \Phi_n,
\]
define
\[
G_\ell := \sum_{n \in I_\ell \cap \{1,\dots,D\}} a_n \Phi_n.
\]
Since $I_0 \cap \{1,\dots,D\} = \emptyset$, we have $G_0=0$ by definition,
and
\[
G = \sum_{\ell \ge 0} G_\ell.
\]

\begin{lemma}\label{lem:block-annulus}
There exist absolute constants $C,c>0$ such that for all $j,\ell \ge 0$,
\[
\frac1\pi \int_{A_j} \abs{G_\ell(z)}^2 \ee^{-\abs{z}^2}\,\dd m(z)
\le
C \exp\!\Bigl(-c\bigl(\abs{j-\ell}-4\bigr)_+^2\Bigr)\norm{G_\ell}^2.
\]
\end{lemma}

\begin{proof}
For each fixed $r$, the functions $\Phi_n(r\cdot)$ with $n \in I_\ell$ have
distinct angular frequencies $n-1$, so they are orthogonal on the circle.
Therefore
\[
\int_{S^1} \abs{G_\ell(r\zeta)}^2\,\dd \sigma(\zeta)
=
\sum_{n \in I_\ell \cap \{1,\dots,D\}}
\abs{a_n}^2 \frac{r^{2n-2}(r^2-n)^2}{n!}.
\]
Integrating over $A_j$ and applying Lemma~\ref{lem:annulus-basis} termwise
gives
\[
\frac1\pi \int_{A_j} \abs{G_\ell(z)}^2 \ee^{-\abs{z}^2}\,\dd m(z)
\le
\sum_{n \in I_\ell \cap \{1,\dots,D\}}
\abs{a_n}^2 C \exp\!\Bigl(-c\bigl(\abs{j-\sqrt n}-3\bigr)_+^2\Bigr).
\]
If $n \in I_\ell$, then $\abs{\sqrt n - \ell}<1$, hence
\[
\bigl(\abs{j-\sqrt n}-3\bigr)_+
\ge
\bigl(\abs{j-\ell}-4\bigr)_+.
\]
Factoring out the common exponential and using
Corollary~\ref{cor:expansion} gives the result.
\end{proof}

\section{Leakage of distant blocks}

Fix an integer $M \ge 1$. For each $j \ge 0$, define
\[
V_j := \sum_{\abs{\ell-j}\le M} G_\ell,
\qquad
R_j := G - V_j.
\]
The function $V_j$ is the local part of $G$ relevant to the annulus $A_j$,
while $R_j$ is the contribution of distant blocks.

\begin{lemma}\label{lem:leakage}
Define
\[
\eta_M
:=
2C \sum_{m>M} \exp\!\bigl(-c(m-4)_+^2\bigr),
\]
where $C,c$ are the constants from Lemma~\ref{lem:block-annulus}. Then
\[
\sum_{j\ge0}
\frac1\pi \int_{A_j} \abs{R_j(z)}^2 \ee^{-\abs{z}^2}\,\dd m(z)
\le
\eta_M \norm{G}^2.
\]
Moreover, $\eta_M \downarrow 0$ as $M \to \infty$.
\end{lemma}

\begin{proof}
For each fixed radius, disjoint blocks are orthogonal on the circle, so on
every annulus
\[
\frac1\pi \int_{A_j} \abs{R_j(z)}^2 \ee^{-\abs{z}^2}\,\dd m(z)
=
\sum_{\abs{\ell-j}>M}
\frac1\pi \int_{A_j} \abs{G_\ell(z)}^2 \ee^{-\abs{z}^2}\,\dd m(z).
\]
Applying Lemma~\ref{lem:block-annulus}, summing in $j$, and swapping the order
of summation yields
\[
\sum_{j\ge0}
\frac1\pi \int_{A_j} \abs{R_j(z)}^2 \ee^{-\abs{z}^2}\,\dd m(z)
\le
\sum_{\ell\ge0}
\sum_{\abs{j-\ell}>M}
C \exp\!\bigl(-c(\abs{j-\ell}-4)_+^2\bigr)\norm{G_\ell}^2.
\]
The inner sum is bounded by
$2C\sum_{m>M}\exp(-c(m-4)_+^2)=\eta_M$, so
\[
\sum_{j\ge0}
\frac1\pi \int_{A_j} \abs{R_j(z)}^2 \ee^{-\abs{z}^2}\,\dd m(z)
\le
\eta_M \sum_{\ell\ge0}\norm{G_\ell}^2
=
\eta_M \norm{G}^2
\]
by Corollary~\ref{cor:expansion}. The Gaussian tail clearly tends to $0$ as
$M\to\infty$.
\end{proof}

\section{The annulus-circle estimate}

The next step is to estimate $V_j$ on each annulus. The key point is that
after dividing by $\Phi_0(re^{it}) = r e^{-it}$, the normalised perturbation
still has only positive frequencies.

\begin{lemma}\label{lem:annulus-circle}
For every integer $M \ge 1$ there exists $B_M>0$ with
\[
B_M^2 \le 10^7 M^2
\]
such that for every $j \ge 0$ and every $r \in [j,j+1]$,
\[
\int_{S^1} \abs{V_j(r\zeta)}^2\,\dd \sigma(\zeta)
\le
B_M^2 \int_{S^1}
\bigl(\abs{\Phi_0 + V_j}(r\zeta) - \abs{\Phi_0}(r\zeta)\bigr)^2
\dd \sigma(\zeta).
\]
\end{lemma}

\begin{proof}
If $r=0$, then necessarily $j=0$. Since $\Phi_0(0)=0$,
\[
\int_{S^1} \abs{V_0(0)}^2\,\dd \sigma
\le
B_M^2 \int_{S^1}
\bigl(\abs{\Phi_0 + V_0}(0) - \abs{\Phi_0}(0)\bigr)^2
\dd \sigma
\]
for every $B_M \ge 1$, because the right-hand integrand is
$\abs{V_0(0)}^2$. Thus it remains to consider $r>0$.

Assume $r>0$. By Lemma~\ref{lem:circle-reduction},
\[
V_j(re^{it}) = r e^{-it} Q_{j,r}(e^{it}),
\]
where $Q_{j,r}$ is a trigonometric polynomial with only strictly positive
frequencies. Moreover,
\[
\abs{V_j(re^{it})} = r\abs{Q_{j,r}(e^{it})},
\qquad
\bigl(\abs{\Phi_0 + V_j} - \abs{\Phi_0}\bigr)(re^{it})
=
r\,\rhoFn\bigl(Q_{j,r}(e^{it})\bigr).
\]
So it is enough to prove
\[
\int_{S^1} \abs{Q_{j,r}}^2\,\dd \sigma
\le
B_M^2 \int_{S^1} \rhoFn(Q_{j,r})^2\,\dd \sigma.
\]

\emph{Low annuli: $j < 250M$.}
If $\abs{\ell-j}\le M$ and $n \in I_\ell$, then
$n < (j+M+1)^2$. Hence $Q_{j,r}$ has at most $(j+M+1)^2$ frequencies. In
this range,
\[
(j+M+1)^2 \le (251M+1)^2 \le 252^2 M^2.
\]
Proposition~\ref{prop:local-circle} therefore gives
\[
\int_{S^1} \abs{Q_{j,r}}^2\,\dd \sigma
\le
144 \cdot 252^2 M^2 \int_{S^1} \rhoFn(Q_{j,r})^2\,\dd \sigma
<
10^7 M^2 \int_{S^1} \rhoFn(Q_{j,r})^2\,\dd \sigma.
\]

\emph{High annuli: $j \ge 250M$.}
Here $j \ge M+1$, so the frequency support of $Q_{j,r}$ is contained in
\[
[(j-M)^2,(j+M+1)^2).
\]
Thus the lowest frequency is
\[
N_{j,M} := (j-M)^2
\]
and the bandwidth is at most
\[
L_{j,M}
:=
(j+M+1)^2 - (j-M)^2
=
(2M+1)(2j+1).
\]
For fixed $M$, the function
\[
x \mapsto \frac{2x+1}{(x-M)^2}
\]
is decreasing on $(M,\infty)$, since
\[
\frac{\dd}{\dd x}\frac{2x+1}{(x-M)^2}
=
-\frac{2(x+M+1)}{(x-M)^3}
<
0.
\]
Therefore $L_{j,M}/N_{j,M}$ is maximal at $j=250M$ among all $j \ge 250M$. At
that point,
\[
L_{250M,M}
=
(2M+1)(500M+1)
\le
1503 M^2,
\qquad
N_{250M,M} = (249M)^2.
\]
Since
\[
1343 \cdot 1503^2 < 249^4,
\]
we obtain
\[
1343\,L_{j,M}^2 \le N_{j,M}^2
\qquad
(j \ge 250M).
\]
Proposition~\ref{prop:high-circle} then yields
\[
\int_{S^1} \abs{Q_{j,r}}^2\,\dd \sigma
\le
32 \int_{S^1} \rhoFn(Q_{j,r})^2\,\dd \sigma
\le
10^7 M^2 \int_{S^1} \rhoFn(Q_{j,r})^2\,\dd \sigma.
\]

Taking $B_M^2 = 10^7 M^2$ proves the lemma.
\end{proof}

\section{Assembly of the argument}

\begin{lemma}\label{lem:assembly-circle}
For every $j \ge 0$ and every $r \in [j,j+1]$,
\begin{align*}
\int_{S^1} \abs{G(r\zeta)}^2\,\dd \sigma(\zeta)
&\le
4B_M^2 \int_{S^1}
\bigl(\abs{\Phi_0 + G}(r\zeta) - \abs{\Phi_0}(r\zeta)\bigr)^2
\dd \sigma(\zeta) \\
&\qquad +
\bigl(4B_M^2 + 2\bigr)\int_{S^1} \abs{R_j(r\zeta)}^2\,\dd \sigma(\zeta).
\end{align*}
\end{lemma}

\begin{proof}
Lemma~\ref{lem:annulus-circle} gives
\[
\int_{S^1} \abs{V_j(r\zeta)}^2\,\dd \sigma
\le
B_M^2 \int_{S^1}
\bigl(\abs{\Phi_0 + V_j}(r\zeta) - \abs{\Phi_0}(r\zeta)\bigr)^2
\dd \sigma.
\]
Since $G = V_j + R_j$, the reverse triangle inequality gives
\[
\Bigl\lvert
\bigl(\abs{\Phi_0 + V_j} - \abs{\Phi_0}\bigr)
-
\bigl(\abs{\Phi_0 + G} - \abs{\Phi_0}\bigr)
\Bigr\rvert
\le
\abs{R_j}.
\]
Hence
\[
\abs{\abs{\Phi_0 + V_j} - \abs{\Phi_0}}
\le
\abs{\abs{\Phi_0 + G} - \abs{\Phi_0}} + \abs{R_j}.
\]
Squaring and using $(a+b)^2 \le 2a^2 + 2b^2$ yields
\begin{align*}
\int_{S^1} \abs{V_j(r\zeta)}^2\,\dd \sigma
&\le
2B_M^2 \int_{S^1}
\bigl(\abs{\Phi_0 + G}(r\zeta) - \abs{\Phi_0}(r\zeta)\bigr)^2
\dd \sigma \\
&\qquad +
2B_M^2 \int_{S^1} \abs{R_j(r\zeta)}^2\,\dd \sigma.
\end{align*}
Finally,
\[
\abs{G}^2 \le 2\abs{V_j}^2 + 2\abs{R_j}^2,
\]
and substituting the previous estimate gives the claim.
\end{proof}

\begin{proof}[Proof of Theorem~\ref{thm:main-orthog}]
We first treat finite sums
\[
G = \sum_{n=1}^D a_n \Phi_n.
\]
Fix an integer $M \ge 1$ for the moment. Multiply the inequality from
Lemma~\ref{lem:assembly-circle} by $2r\ee^{-r^2}$, integrate over
$r \in [j,j+1]$, and sum in $j \ge 0$. Since the annuli $A_j$ partition the
plane, we obtain
\[
\norm{G}^2
\le
4B_M^2 \norm{\abs{\Phi_0 + G} - \abs{\Phi_0}}^2
+
\bigl(4B_M^2 + 2\bigr)
\sum_{j\ge0}
\frac1\pi \int_{A_j} \abs{R_j(z)}^2 \ee^{-\abs{z}^2}\,\dd m(z).
\]
By Lemma~\ref{lem:leakage},
\[
\norm{G}^2
\le
4B_M^2 \norm{\abs{\Phi_0 + G} - \abs{\Phi_0}}^2
+
\bigl(4B_M^2 + 2\bigr)\eta_M \norm{G}^2.
\]
Now $B_M^2 \le 10^7 M^2$, while
\[
\eta_M = 2C \sum_{m>M} \exp\!\bigl(-c(m-4)_+^2\bigr)
\]
has Gaussian decay. Hence
\[
\bigl(4B_M^2 + 2\bigr)\eta_M \longrightarrow 0
\qquad
(M\to\infty).
\]
Choose $M$ so large that
\[
\bigl(4B_M^2 + 2\bigr)\eta_M \le \frac12.
\]
Then
\[
\norm{G}^2
\le
8B_M^2 \norm{\abs{\Phi_0 + G} - \abs{\Phi_0}}^2,
\]
which proves the theorem for finite sums.

For general $G \in H_1 \cap \Phi_0^\perp$, write
\[
G = \sum_{n=1}^\infty a_n \Phi_n
\]
and let
\[
G_N := \sum_{n=1}^N a_n \Phi_n.
\]
By Corollary~\ref{cor:expansion}, $G_N \to G$ in $H_1$. The finite-sum bound
gives
\[
\norm{G_N}
\le
C_0 \norm{\abs{\Phi_0 + G_N} - \abs{\Phi_0}},
\]
with $C_0 = \sqrt{8}\,B_M$ independent of $N$. Since the modulus is
$1$-Lipschitz,
\[
\abs{\bigl(\abs{\Phi_0 + G_N} - \abs{\Phi_0}\bigr)
-
\bigl(\abs{\Phi_0 + G} - \abs{\Phi_0}\bigr)}
\le
\abs{G_N - G},
\]
so the right-hand side converges to
$\norm{\abs{\Phi_0 + G} - \abs{\Phi_0}}$. Passing to the limit yields the
theorem.
\end{proof}

\begin{remark}
The only new analytic ingredient is Lemma~\ref{lem:annulus-basis}. Once the
annulus localisation is available, the remaining argument follows the same
block decomposition and leakage scheme as in the analytic case, provided one
also keeps track of the explicit $M$-dependence of $B_M$, passes from the
phase-normalised orthogonal reduction to phase-optimised stability via
Lemma~\ref{lem:phase-alignment}, and squares the signed defect only after
inserting absolute values as in Lemma~\ref{lem:assembly-circle}.
\end{remark}

\begin{thebibliography}{1}

\bibitem{FockSPR}
J.~de~Dios~Pont and M.~Taylor,
\emph{Local Stable Phase Retrieval in Fock Space},
Lean-verified repository version \texttt{FOCK Old/FockSPR-10.pdf}, 2026.

\end{thebibliography}

\end{document}
